При фиксированных локальных коэффициентах (\hat f'_j) обновление направления (\beta) записывается как ridge-задача

[
\hat\beta
=========

\arg\min_{\beta\in\mathbb R^d}
\left{
\sum_{j\in\mathcal J}
\left|
I_j-\hat f'_jU_j\beta
\right|*2^2
+
\lambda|\beta-\beta*{\mathrm{init}}|_2^2
\right}.
]

Определим блочный оператор

[
A=
\begin{pmatrix}
\hat f'_1U_1\
\vdots\
\hat f'_J U_J
\end{pmatrix},
\qquad
b=
\begin{pmatrix}
I_1\
\vdots\
I_J
\end{pmatrix}.
]

Тогда задача эквивалентна обычной задаче наименьших квадратов

[
\hat\beta
=========

\arg\min_\beta
\left|
\begin{pmatrix}
A\
\sqrt{\lambda}I_d
\end{pmatrix}
\beta
-----

\begin{pmatrix}
b\
\sqrt{\lambda}\beta_{\mathrm{init}}
\end{pmatrix}
\right|_2^2.
]

Она решается методом LSMR без построения матрицы (A). Достаточно реализовать операции

[
Av
==

\operatorname{col}_{j\in\mathcal J}
\left(\hat f'_jU_jv\right),
]

[
A^\top u
========

\sum_{j\in\mathcal J}
\hat f'_jU_j^\top u_j.
]

Для расширенного ridge-оператора:

[
\widetilde A v
==============

\begin{pmatrix}
Av\
\sqrt{\lambda}v
\end{pmatrix},
\qquad
\widetilde A^\top
\begin{pmatrix}
u\
z
\end{pmatrix}
=============

A^\top u+\sqrt{\lambda}z.
]

LSMR решает исходную least-squares задачу, не формируя нормальные уравнения

[
(A^\top A+\lambda I)\beta
=========================

A^\top b+\lambda\beta_{\mathrm{init}}.
]

Поэтому обусловленность не возводится в квадрат, как при решении через CG по (A^\top A). Стоимость одной итерации составляет

[
O!\left(Jn_\Phi d\right),
]

а дополнительная память не зависит от размера полной блочной матрицы. После решения направление нормируется:

[
\beta\leftarrow\frac{\beta}{|\beta|_2}.
]

Полученное (\beta) фиксируется, после чего обновляются локальные коэффициенты и начинается следующая итерация чередующейся оптимизации.

